\section{Introduction: why soil water moves}
The previous chapter described how much water a soil holds and how tightly. This chapter describes how the water moves. Movement is the reason that an irrigation applied at the surface wets the root zone, that a heavy rain drains below it, that a dry layer draws water from a wet one and that a shallow water table feeds a crop. It is also the reason that the roots of a plant meet a supply that decays as the soil dries. The equation that governs all this was written by Richards in 1931 by joining two earlier results: Darcy's law, which says how fast water moves under a gradient of head, and the principle of conservation of mass, which says that what flows into a volume and does not flow out must stay in it.

A textbook of soil physics divides the subject by the state of the soil (saturated or unsaturated), by the time behaviour (steady or transient) and by the dimension (one, two or three). This book treats the \emph{one-dimensional, unsaturated, transient} case, which is the case of Daisy and of the application, because a root zone of a uniform field receives its water from above and loses it below. Sections of this chapter give the types of flow, the derivation of the equation, its boundary conditions, the sink term of the roots, the numerical solution and the structure of the Daisy model that implements them.

\section{Types of flow in a soil}
Flow is called \kw{saturated} when all pores are filled and the matric head is zero or positive; the conductivity is then the constant $K_s$ and Darcy's law applies in its simplest form. Flow is \kw{unsaturated} when air occupies part of the pores; the conductivity $K(\theta)$ falls with the water content by several orders of magnitude. Flow is \kw{steady} when the flux at each point does not change with time (a constant evaporation from a fixed water table is the standard case) and \kw{transient} otherwise (infiltration, redistribution, drying). Three further divisions are used in practice: \emph{matrix flow}, through the fine pores that Richards' equation describes; \emph{preferential flow}, through cracks and root channels that bypass the matrix; and \emph{vapour flow}, which matters only in the top centimetres of a dry soil. Daisy has a matrix domain and a macropore domain; the one-dimensional matrix flow is the part that this book works with.

\section{Darcy's law and the continuity equation}
Darcy (1856) found that the volume flux $q$ of water through a saturated sand column is proportional to the difference of hydraulic head $H$ along the column and inversely proportional to its length. Written for the unsaturated soil by Buckingham (1907), with $z$ positive upward, it is

\begin{equation}q=-K(h)\,\frac{\partial H}{\partial z}=-K(h)\Bigl(\frac{\partial h}{\partial z}+1\Bigr),\label{eq:darcy}\end{equation}


where $q$ is positive upward, $K(h)$ is the unsaturated hydraulic conductivity and the $+1$ is the gradient of the gravitational head. Conservation of mass for a thin slice, with the sink $S$ of root uptake (volume of water per unit volume of soil per unit time), is

\begin{equation}\frac{\partial\theta}{\partial t}=-\frac{\partial q}{\partial z}-S(z,t).\label{eq:cont}\end{equation}
\par\noindent where $\theta$ is the water content, $t$ the time (d), $q$ the flux (cm\,d$^{-1}$), $z$ the depth (cm) and $S(z,t)$ the sink of root uptake (d$^{-1}$).



Substituting (\ref{eq:darcy}) in (\ref{eq:cont}) and writing $C(h)=\mathrm d\theta/\mathrm dh$ for the specific water capacity gives the Richards equation in its head form

\begin{equation}C(h)\frac{\partial h}{\partial t}=\frac{\partial}{\partial z}\Bigl[K(h)\Bigl(\frac{\partial h}{\partial z}+1\Bigr)\Bigr]-S(z,t).\label{eq:richards}\end{equation}
\par\noindent where $C(h)$ is the water capacity (cm$^{-1}$), $h$ the suction (cm), $t$ the time (d), $z$ the depth (cm), $K(h)$ the conductivity (cm\,d$^{-1}$) and $S(z,t)$ the root uptake (d$^{-1}$).



\exbox{Darcy flux in a drying column}{In a column of loam, the suction is 80 cm at 20 cm depth and 120 cm at 40 cm depth (z positive upward, surface at 0). The conductivity at the mean suction of 100 cm is 0.65 cm\,d$^{-1}$. Find the heads, the gradient and the flux, and say whether the soil is draining or wetting at this depth.}\sol With $z$ upward, $z_{20}=-20$ and $z_{40}=-40$. Heads: $H_{20}=-80-20=-100$ cm, $H_{40}=-120-40=-160$ cm. The gradient \[\partial H/\partial z=(H_{20}-H_{40})/(z_{20}-z_{40})=(-100+160)/(20)=3.0\] (dimensionless). Flux \[q=-K\,\partial H/\partial z=-0.65\times3.0=-1.95\ \text{cm\,d$^{-1}$}\] \ans The sign is negative, so the flow is \emph{downward}: the column is draining at 19.5 mm per day, mostly under the gradient of suction, not gravity alone. \finans

\section{Boundary and initial conditions}
Equation~(\ref{eq:richards}) needs one condition at each end of the column and one initial profile. Three types of boundary condition are used. A \kw{Dirichlet} condition fixes the head (a ponded surface, $h=0$ or a water table, $h=0$ at its depth). A \kw{Neumann} condition fixes the flux (a known rain or evaporation rate, or zero flux across an impermeable layer). A \kw{free drainage} condition sets the gradient of head to zero at the lower boundary, so that $q=-K$: water leaves at the rate of the conductivity. At the surface, Daisy and the application switch between the Neumann condition (the rain or irrigation rate) and the Dirichlet condition (ponding) according to which one the soil can accept: if the rate exceeds the infiltration capacity, the surplus becomes surface storage or runoff.

\section{Water uptake by roots}
Roots take water from the soil to supply transpiration. Two families of models describe this. In \kw{microscopic} models each root is a cylinder of radius $r_0$ in a soil cylinder of radius $r_1$, and the radial flow obeys the Richards equation in cylindrical form; the flux depends on the difference between the suction at $r_1$ and at the root surface. In \kw{macroscopic} models the whole root system is a distributed sink $S(z)$. Feddes' form multiplies the potential sink $S_p(z)$, a share of the potential transpiration given by the root density, by a reduction factor $\alpha(h)$ that is one in the wet range, falls linearly to zero between a critical suction $h_3$ (and $h_4$, the wilting point), and is zero in a flooded soil (below $h_1$, $h_2$). Daisy follows the microscopic idea: the sink in each layer is the lesser of a demand and the radial conductance of the soil around the roots times the head difference, with the root density as the length of root per unit volume. The conductance of the path depends on the conductivity of the soil, so that uptake falls as the soil dries, which is why the reach of a suction trigger (Chapter~4) depends on $K$.

\begin{equation}S(z)=\alpha(h)\,S_p(z),\qquad S_p(z)=T_p\,\frac{R(z)}{\int_0^{Z_r}R(z')\,\mathrm dz'},\qquad \alpha(h)=\begin{cases}\dfrac{h-h_4}{h_3-h_4}&h_4<h<h_3\\[4pt]1&h_3\le h\le h_2\\0&\text{otherwise}\end{cases}\label{eq:feddes}\end{equation}
\par\noindent where $S(z)$ is the uptake at depth $z$ (d$^{-1}$), $\alpha(h)$ the Feddes stress function (0 to 1), $S_p(z)$ the potential uptake, $T_p$ the potential transpiration (cm\,d$^{-1}$), $R(z)$ the root density at depth $z$, $Z_r$ the rooting depth, $z'$ the variable of integration and $h_2,h_3,h_4$ the suction limits of the stress function (cm).

\exbox{Feddes uptake in four layers}{A root zone has four layers with relative root densities 0.40, 0.30, 0.20, 0.10 (top to bottom) and suctions $|h|$ of 150, 500, 2000 and 8000 cm. The potential transpiration is $T_p=5.0$ mm\,d$^{-1}$; the critical suction is $|h_3|=400$ cm and the wilting suction is 15\,849 cm. Find the actual uptake of each layer and the total.}\sol \givenline{$T_p=5.0$ mm\,d$^{-1}$; ; $|h_3|=400$ cm}\ans The potential sinks are \[S_p=T_p R_i/\sum R=5.0\times(0.40,0.30,0.20,0.10)=(2.00,1.50,1.00,0.50)\ \text{mm\,d$^{-1}$}\]
The reduction factors: layer 1 ($|h|=150<400$): $\alpha=1$; layer 2 ($500$): \[\alpha=(15849-500)/(15849-400)=0.994\]
layer 3 ($2000$): \[\alpha=(15849-2000)/15449=0.896\]
layer 4 ($8000$): \[\alpha=(15849-8000)/15449=0.508\]
Actual sinks: $2.00$, $1.49$, $0.90$, $0.25$ mm\,d$^{-1}$; total $=4.64$ mm\,d$^{-1}$ \ans The roots in the wet top layer supply 2.0 mm, and the deficit of 0.36 mm is not made good by the lower layers because they are drier: in Feddes' simple form it is a stress; Daisy lets the wetter layers compensate to a degree set by the root conductance. \finans

\section{Numerical solution of the Richards equation}
The Richards equation has no closed-form solution for a general soil, so it is solved numerically. The profile is divided into nodes at spacing $\Delta z$, time into steps $\Delta t$, and derivatives are replaced by differences. In an \kw{explicit} scheme the new values are computed from the old ones only; it is simple but stable only if $\Delta t\le \Delta z^2/(2D_{\max})$, where $D=K/C$ is the soil-water diffusivity, which reaches hundreds of cm$^2$\,d$^{-1}$ in a wet sand and forces tiny steps. In an \kw{implicit} scheme the unknowns at the new time are found together by solving a tridiagonal system, at the cost of iteration for the non-linearity of $K$ and $C$; it is unconditionally stable and is what Daisy and most models use, with a time step that shrinks when the iteration does not converge (typically from 1 h to a few seconds during a rain). A well-made scheme conserves mass: the change of storage equals the net flux across the boundaries minus the uptake, and the mass-balance error is reported and kept below 0.1 mm.

\exbox{Stability limit of the explicit scheme}{For the loam S06 ($K_s=25$ cm\,d$^{-1}$ assumed, $\alpha=0.036$, $n=1.56$, $\theta_s-\theta_r=0.448$) find the diffusivity $D=K/C$ at suctions of 50, 100, 300 and 1000 cm and the largest time step of an explicit scheme with $\Delta z=2$ cm.}\sol \givenline{$K_s=25$ cm\,d$^{-1}$ assumed ; $\alpha=0.036$ ; $n=1.56$ ; $\theta_s-\theta_r=0.448$ ; $\Delta z=2$ cm}\ans The specific capacity is \[C=(\theta_s-\theta_r)\,\alpha n m(\alpha h)^{n-1}[1+(\alpha h)^n]^{-m-1}\] with $m=0.359$, and $K=K_s K_r$. at $|h|=50$ cm: $K=0.258$ cm\,d$^{-1}$, $C=0.00229$ cm$^{-1}$, $D=112.9$ cm$^2$\,d$^{-1}$, $\Delta t_{max}=0.0177$ d (25.5 min); at $|h|=100$ cm: $K=0.034$ cm\,d$^{-1}$, $C=0.00103$ cm$^{-1}$, $D=33.0$ cm$^2$\,d$^{-1}$, $\Delta t_{max}=0.0606$ d (87.3 min); at $|h|=300$ cm: $K=0.001$ cm\,d$^{-1}$, $C=0.00021$ cm$^{-1}$, $D=4.5$ cm$^2$\,d$^{-1}$, $\Delta t_{max}=0.4489$ d (646.4 min); at $|h|=1000$ cm: $K=0.000$ cm\,d$^{-1}$, $C=0.00003$ cm$^{-1}$, $D=0.5$ cm$^2$\,d$^{-1}$, $\Delta t_{max}=4.0984$ d (5901.8 min) \ans The explicit step is a few minutes in the wet range and the 1-h step of Daisy would be unstable, which is why an implicit scheme is used. \finans

\section{The structure of the Daisy model}
Daisy is a one-dimensional soil--plant--atmosphere model developed at the University of Copenhagen since the 1980s (Hansen et al., 2012). It is organised in \emph{modules} that exchange state each time step: the \kw{weather} module supplies rain, temperature, radiation, humidity and wind; the \kw{surface} module partitions rain into interception, ponding, runoff and infiltration and computes evaporation; the \kw{soil water} module solves Richards' equation (and, if requested, a bypass flow in macropores); the \kw{soil heat} module solves the heat equation for the soil temperature; the \kw{solute} and \kw{organic matter} modules transport nitrate and mineralise organic nitrogen; the \kw{vegetation} module grows a crop with a canopy, roots and a yield; and the \kw{management} module applies tillage, sowing, fertiliser, irrigation and harvest. Each module is a family of alternative \emph{models} selected in a description file; the experiments of this book change a handful of these choices (soil, trigger, mulch, tillage, climate) over a fixed base case and read the water balance of the result.

A Daisy run is described by a text file of nested entries. The user, or the application on the user's behalf, writes the file, runs the program for the chosen years and reads the log files (daily and yearly) that the program writes. The application never changes the physics of Daisy; it \emph{reads} the archived results of 10\,190 simulations and, for the layers that need a model, uses an engine that reproduces the same relations in a few lines (Chapter~9).

